\chapter*{Glosario de siglas y términos}
\addcontentsline{toc}{chapter}{Glosario}

\begin{description}

\item[APT] \textit{Advanced Persistent Threat}. Actor con capacidad, recursos y motivación
para mantener una intrusión prolongada y sigilosa en una organización objetivo.

\item[ATT\&CK] Base de conocimiento pública de MITRE que cataloga tácticas, técnicas y
procedimientos (TTP) empleados por adversarios reales, organizados en una matriz.

\item[CEP] \textit{Complex Event Processing}. Paradigma de procesamiento que detecta
patrones formados por secuencias de eventos relacionados en el tiempo, en lugar de evaluar
eventos de forma aislada.

\item[CTI] \textit{Cyber Threat Intelligence}. Inteligencia de amenazas: información
contextualizada sobre actores, campañas, herramientas e indicadores de compromiso.

\item[CVE] \textit{Common Vulnerabilities and Exposures}. Identificador público y único
asignado a una vulnerabilidad conocida.

\item[CVSS] \textit{Common Vulnerability Scoring System}. Sistema estándar de puntuación de
la severidad de una vulnerabilidad.

\item[Datalog] Lenguaje declarativo de consulta y razonamiento, subconjunto de Prolog,
empleado para derivar hechos nuevos a partir de hechos base y reglas de inferencia.

\item[DLQ] \textit{Dead Letter Queue}. Cola de rechazos donde se derivan los mensajes que
no han podido procesarse, para su inspección posterior sin pérdida de información.

\item[EPSS] \textit{Exploit Prediction Scoring System}. Estimación de la probabilidad de
que una vulnerabilidad concreta sea explotada en un plazo determinado.

\item[Event time] Marca temporal correspondiente al instante en que ocurrió realmente un
evento, por oposición al \textit{processing time} (instante en que el sistema lo procesa).

\item[IOC] \textit{Indicator of Compromise}. Artefacto observable que evidencia una
intrusión: una dirección IP, un hash, un dominio o una ruta de fichero.

\item[Kill chain] Modelo que describe un ataque como una secuencia ordenada de fases, desde
el reconocimiento inicial hasta la consecución del objetivo.

\item[LSASS] \textit{Local Security Authority Subsystem Service}. Proceso de Windows que
gestiona la autenticación y custodia credenciales en memoria; objetivo habitual de las
técnicas de robo de credenciales.

\item[MulVAL] \textit{Multihost, multistage Vulnerability Analysis}. Analizador de
seguridad de red basado en lógica que emplea Datalog para derivar rutas de ataque
factibles a partir de la configuración y las vulnerabilidades de una red.

\item[OCSF] \textit{Open Cybersecurity Schema Framework}. Esquema abierto y agnóstico de
proveedor para la normalización de eventos de seguridad.

\item[Sigma] Formato abierto y agnóstico para expresar reglas de detección sobre registros
de eventos, convertible a los lenguajes de consulta de distintos SIEM.

\item[SIEM] \textit{Security Information and Event Management}. Plataforma de recolección,
correlación y análisis centralizado de eventos de seguridad.

\item[SOC] \textit{Security Operations Centre}. Centro de operaciones de seguridad: equipo
responsable de la monitorización, detección y respuesta ante incidentes.

\item[STIX] \textit{Structured Threat Information eXpression}. Lenguaje estandarizado por
OASIS para representar e intercambiar inteligencia de amenazas de forma estructurada.

\item[Sysmon] \textit{System Monitor}. Herramienta de Microsoft que amplía la auditoría
nativa de Windows con telemetría detallada de procesos, red, ficheros y registro.

\item[TTP] \textit{Tactics, Techniques and Procedures}. Descripción del comportamiento de un
adversario en tres niveles de abstracción: el objetivo perseguido, el método empleado y su
implementación concreta.

\item[Watermark] Mecanismo de los sistemas de procesamiento en \textit{streaming} que estima
el progreso del tiempo de evento, permitiendo decidir cuándo una ventana temporal puede
considerarse completa.

\end{description}
